\setcounter{aufgabe}{37}
\hypertarget{e5}{}\einheitenkopf{Einheit 5 von 5 · Sachaufgaben}

\begin{aufgabe}{Sachaufgaben -- Welche Lösung passt? Die Rechnung liefert zwei
Lösungen. Kreuze an, welche zur Frage passt. Rechne nichts.}
\begin{teile}
\teil Gesucht ist die Seitenlänge eines Quadrats. Lösungen: $7$ und $-7$. \\
\kreuz{nur die positive} \kreuz{nur die negative} \kreuz{beide}
\teil Gesucht sind alle Zahlen, deren Quadrat $64$ ist. Lösungen: $8$ und $-8$. \\
\kreuz{nur die positive} \kreuz{nur die negative} \kreuz{beide}
\teil Gesucht ist die Anzahl der Schüler. Lösungen: $12$ und $-5$. \\
\kreuz{nur die positive} \kreuz{nur die negative} \kreuz{beide}
\teil Gesucht ist die Breite eines Rechtecks in Metern. Lösungen: $-9$ und $4$. \\
\kreuz{nur die positive} \kreuz{nur die negative} \kreuz{beide}
\teil Gesucht sind alle Temperaturen in $^\circ$C, die eine Bedingung erfüllen.
Lösungen: $-3$ und $6$. \\
\kreuz{nur die positive} \kreuz{nur die negative} \kreuz{beide}
\end{teile}
\end{aufgabe}

\begin{aufgabe}{Sachaufgaben -- Zahlenrätsel. Stell die Gleichung auf, löse sie und
gib alle Zahlen an, die passen.}
\beispiel{\text{Das Quadrat einer Zahl ist } 36. \\
x^2 &= 36 \\ x_1 &= 6 && x_2 = -6}
\begin{teile}
\teil Das Quadrat einer Zahl ist $16$.
\teil Das Quadrat einer Zahl ist $196$.
\teil Das Quadrat einer Zahl ist $400$.
\teil Das Quadrat einer Zahl ist $2{,}25$.
\end{teile}
Hier hilft kein reines Wurzelziehen mehr. Stell die Gleichung auf, ordne und löse.
\begin{teile}
\setcounter{teil}{4}
\teil Das Quadrat einer Zahl ist um $24$ größer als das Fünffache der Zahl.
\teil Das Produkt einer Zahl mit ihrem Nachfolger ist $56$.
\teil Das Produkt einer Zahl mit der um $4$ größeren Zahl ist $45$.
\end{teile}
\end{aufgabe}

\begin{aufgabe}{Sachaufgaben -- Rechteck. Nenn zuerst die Unbekannte ($x = \ldots$),
stell die Gleichung auf, löse sie, schließe die unpassende Lösung aus und antworte
mit Einheit.}
\begin{teile}
\teil Ein rechteckiges Beet ist doppelt so lang wie breit und hat einen
Flächeninhalt von $72\,\text{m}^2$. Berechne Breite und Länge.
\teil Ein Rechteck ist um $5\,$m länger als breit und hat einen Flächeninhalt von
$84\,\text{m}^2$. Berechne Breite und Länge. Die Skizze zeigt die Bezeichnungen.
\end{teile}

\rechteck[punkte=,seiten={x + 5,x,,}]{5}{3}

\begin{teile}
\setcounter{teil}{2}
\teil Ein rechteckiges Grundstück hat einen Umfang von $34\,$m und einen
Flächeninhalt von $60\,\text{m}^2$. Berechne die beiden Seitenlängen.
\end{teile}
\end{aufgabe}

\begin{aufgabe}{Sachaufgaben -- Fehler finden. Ein Rechteck ist um $7\,$m länger als
breit und hat einen Flächeninhalt von $60\,\text{m}^2$. Ben rechnet:}
\rechnung{x \cdot (x + 7) &= 60 \\
x^2 + 7x - 60 &= 0 \\
x_1 &= 5 && x_2 = -12}
Er antwortet: „Die Seiten sind $5\,$m und $-12\,$m.`` Finde den Fehler, benenne ihn
und gib die richtige Antwort.

\begin{teile}
\teil Ein Rechteck ist um $2\,$m länger als breit und hat einen Flächeninhalt von
$48\,\text{m}^2$. Berechne die Seiten und antworte mit Einheit.
\end{teile}
\end{aufgabe}

\begin{aufgabe}{Sachaufgaben -- Begründen. Antworte in ganzen Sätzen.}
\begin{teile}
\teil Warum darf beim Rechteck nur eine der beiden Lösungen als Seitenlänge gelten,
beim Zahlenrätsel dagegen oft beide? Begründe.
\teil Eine Sachaufgabe führt auf die Gleichung $x^2 + 6x + 20 = 0$. Was bedeutet es
für die Aufgabe, dass diese Gleichung keine Lösung hat? Begründe.
\end{teile}
\end{aufgabe}
